\baselineskip12.2pt plus.2pt minus.2pt
\newcommand{\ppmod}[1]{~(\textrm{mod #1})} % alternatywny pmod, bliżej równości


\spis{Świat jest pozytywny i~kwadratowy}
{Jacek Pytlak, Mariusz Skałba}

\wtyt{Świat jest pozytywny i~kwadratowy}

\waut{Jacek PYTLAK, Mariusz SKAŁBA*}
\aafil{Wydział Matematyki, Informatyki i~Mechaniki, Uniwersytet Warszawski}

Gdy ktoś zaczepi Was na ulicy i~zapyta: \textit{Co daje minus razy minus?,} to w~zależności od nastroju odpowiecie bardzo krótko: {\it  To nic nie daje} albo też krótko: {\it To daje bardzo dodatni plus. }
Dlatego najpierw przypomnimy fragmenty  teorii ogólnej, niezależnej od
koniunktury.

Niech $p$ będzie nieparzystą liczbą pierwszą i~niech
$\mathbb{F}_p=\{0,\ldots,p-1\}$ będzie zbiorem reszt z~dzielenia przez
$p$. Wówczas
dla niektórych reszt $r\neq 0$ istnieją dwie reszty $x$ spełniające równanie
$$
x^2\equiv r\pmod{p},
$$ 
a~dla pozostałych niezerowych $r$ nie ma w~ogóle takich
reszt $x$. W~pierwszym przypadku mówimy, że
$r$ jest {\it resztą kwadratową modulo $p$}, a~w~drugim, że $r$ {\it jest
nieresztą kwadratową modulo $p$}. 
I~tak na przykład z~poniższej tabelki dla $p=11$\smallskip
$$\def\arraystretch{1.3}
\definecolor{mygray}{gray}{0}
\begin{tabular}{|c|c|c|c|c|c|c|c|c|c|c|}
\hline
\color{mygray}$x$ &\color{mygray} 1 &\color{mygray} 2 &\color{mygray} 3 &\color{mygray} 4 &\color{mygray} 5 &\color{mygray} 6  &\color{mygray} 7 &\color{mygray} 8 &\color{mygray} 9 &\color{mygray} 10  \\
\hline
\color{magenta}$x^2$ &\color{magenta} 1 &\color{magenta} 4 &\color{magenta} 9 &\color{magenta} 5 &\color{magenta} 3 &\color{magenta} 3 &\color{magenta} 5 &\color{magenta}  9 &\color{magenta} 4  &\color{magenta} 1 \\
\hline
\end{tabular}\smallskip
$$
wynika, że $1,3,4,5,9$ to reszty kwadratowe modulo $11$, a~$2,6,7,8,10$ to
niereszty kwadratowe modulo $11$. To, że wiersz $x^2$ powyższej tabelki jest
środkowo symetryczny, ma charakter ogólny, gdyż zawsze $(p-x)^2\equiv
x^2\pmod{p}$. Zatem reszty z~dzielenia przez $p$ liczb \begin{equation}
\label{sq}
1^2,\,2^2,\,3^2,\,\ldots,\, \left(\frac{p-1}{2}\right)^2\tag{$*$}
\end{equation}
to reszty kwadratowe modulo $p$, a~pozostałe liczby z~ciągu $1,2,3,\ldots, p-1$ to niereszty. 

Jak stwierdzić, czy dana reszta $a$ z~dzielenia przez $p$ jest resztą, czy też
nieresztą kwadratową, nie przeszukując ciągu (\ref{sq})? Mówi o~tym znane kryterium Eulera:
$$
\textrm{$a$ jest resztą
kwadratową modulo $p$}\iff a^{(p-1)/2}\equiv
1\pmod{p}.
$$
 Od tej chwili będziemy zakładać,  że  liczba pierwsza $p$ jest postaci $p=4k+3$.
 Można łatwo udowodnić, że takich liczb pierwszych jest nieskończenie wiele. Trudny jest natomiast dowód asymptotyki dla funkcji 
$\pi_{4,3}(x)$, która (z~definicji) podaje liczbę liczb pierwszych $p$
mniejszych równych $x$, spełniających kongruencję $p\equiv 3\pmod{4}$:
$$
\lim_{x\to\infty}\frac{\pi_{4,3}(x)}{x/\log x}=\frac{1}{2}.
$$
Przypomnijmy, że analogiczna równość zachodzi dla samej funkcji $\pi(x)$, zliczającej liczbę wszystkich
liczb pierwszych nie większych od $x$, przy czym w~granicy po prawej stronie
dostajemy 1.
Statystycznie więc co druga liczba pierwsza jest interesującej nas postaci. 
Zauważmy, że na mocy kryterium Eulera dla
$p=4k+3$ reszta $p-1$ jest nieresztą kwadratową, gdyż
$$
(p-1)^{(p-1)/2}\equiv (-1)^{2k+1}\equiv-1\not\equiv 1\pmod{p}.
$$
Ogólniej: {\it z~dwóch reszt $a$ oraz $p-a$ jedna jest resztą, natomiast druga nieresztą kwadratową.}

Przedstawione obserwacje pozwalają
 nam określić {\it pierwiastek arytmetyczny} z~dowolnej reszty kwadratowej $a$. Niech mianowicie 
$$
b= a^{(p+1)/4} \pmod{p}.
$$
Mamy teraz 
$$
b^2\equiv a^{(p+1)/2}= a^{(p-1)/2}\cdot a\equiv a\pmod{p},
$$
czyli rzeczywiście powyżej określone $b$ jest jednym z~pierwiastków kwadratowych
z~elementu $a$. A~my dodatkowo  uważamy, że $b$ jest tym
\textit{lepszym}
pierwiastkiem i~chcemy go nazywać {\it pierwiastkiem arytmetycznym}!
Przypomnijmy, że w~zakresie liczb rzeczywistych pierwiastkiem kwadratowym
arytmetycznym z~dodatniej liczby $a$ nazywamy taką (jedyną) dodatnią liczbę $b$,
że \vadjust{\goodbreak}$b^2=a$.  
W~zbiorze liczb rzeczywistych $a\in
\mathbb{R}$ jest dodatnie wtedy i~tylko wtedy, gdy $a=b^2$ dla pewnego
$b\in\mathbb R$.
Od tej chwili niech ,,bycie kwadratem'' będzie namiastką
,,dodatniości''.  
U~nas $b$ jako potęga reszty kwadratowej $a$ jest też resztą
kwadratową! Uzasadnia to traktowanie $b$ jako pierwiastka arytmetycznego z~$a$.

Niestety przedstawiona analogia jest ułomna, gdyż w~doskonałym
(czyli rzeczywistym) świecie suma elementów dodatnich jest dodatnia, a~w~świecie
reszt nie musi tak być. Dla przykładu
$$
3\equiv 5^2\ppmod{11}\ \mbox{ oraz }\ 5\equiv 4^2\ppmod{11},\ \mbox{ ale }\ 
8\not\equiv b^2\ppmod{11}\ \mbox{ dla }b\in \mathbb F_{11}. 
$$
Takie ciała jak $\mathbb R$ czy też $\mathbb Q$ wyróżniają się tym, że można w~nich wprowadzić relację porządku, która jest zgodna z~oboma działaniami, tzn. zarówno
iloczyn, jak i~suma elementów dodatnich jest dodatnia. Natomiast ciało liczb
zespolonych oraz ciała skończone nie są tak porządnie uporządkowane. Tak więc
nasz świat rzeczywisty, chociaż tak niedoskonały,  nie jest najgorszym z~możliwych do pomyślenia  światów. 

 Jest też w~użyciu matematycznym bardziej naturalne i~praktyczne pojęcie
 ,,dodatniości'' w~zbiorze reszt niż powyższe ,,bycie resztą
 kwadratową''. Jeśli mianowicie $p$ jest dowolną liczbą pierwszą nieparzystą, to
 zbiór liczb całkowitych w~przedziale $(-p/2,p/2)$ jest zbiorem reprezentantów
 ciała $\mathbb F_p$ reszt z~dzielenia przez $p$. Od teraz klasę reszt mod $p$ będziemy
 nazywać {\it dodatnią}, gdy ma swojego przedstawiciela w~przedziale $(0,p/2)$,
 a~{\it ujemną}, gdy ma przedstawiciela w~$(-p/2,0)$. 

Ważnym zastosowaniem tej dychotomii jest lemat Gaussa: najwygodniejszy portal do
 \marg{Prawo wzajemności reszt kwadratowych orzeka, że jeśli $p$ i $q$ są
 liczbami pierwszymi i~choć jedna z~nich daje resztę~1  z~dzielenia
 przez 4, to $p$ jest resztą kwadratową modulo $q$ wtedy i~tylko wtedy, gdy
 $q$ jest resztą kwadratową modulo $p$. Jeśli obie liczby dają resztę 3 z~dzielenia przez 4, to $p$ jest resztą kwadratową modulo $q$ wtedy i~tylko wtedy, gdy
 $q$ \emph{nie jest} resztą kwadratową modulo $p$.}
Prawa Wzajemności dla Reszt Kwadratowych (PWdRK) -- Gauss użył go w~swoim piątym
dowodzie PWdRK. Pozwolimy tu sobie na dygresję historyczną: przed Gaussem
nikt nie potrafił udowodnić PWdRK. Z~czwórki gigantów arytmetyki wyższej
wymienionych przez Gaussa w~przedmowie do {\it Disquisitiones Arithmeticae}:
Fermata, Eulera, Lagrange'a i~Legendre'a, tylko  ten ostatni był bliski
poprawnego dowodu.  Sam Gauss podał sześć dowodów; po Gaussie opublikowano ponad
250 dowodów! Ta ogromna liczba różnych dowodów słynnego twierdzenia świadczy
pośrednio o~jego wielorakich powiązaniach z~resztą matematyki. 

To wszystko jest bardzo piękne, ale my teraz nie o~tym\dots\ Przy powyższej
definicji  ,,dodatniości''  nie jest prawdą, że nasza metoda
pierwiastkowania w~$\mathbb F_p^{\ast}$ dla ${p\equiv 3\pmod{4}}$ zawsze daje
pierwiastek ,,dodatni'' (tzn. resztę z~przedziału~$(0,p/2)$), ale okazuje się,
że jednak trochę częściej  niż ,,ujemny'': {\it w~przedziale $(0,p/2)$ jest
mianowicie więcej reszt niż niereszt! } Nie ma prostego uzasadnienia tej
prawidłowości, warto jednak wspomnieć, z~jakiego
ogólniejszego i~ważnego twierdzenia to wynika. Dla $p=3$ teza jest oczywista. Załóżmy w~dalszym ciągu, że $p\geq 7$. Niech $\cal R$ oznacza liczbę reszt kwadratowych
modulo~$p$ w~przedziale $(0,p/2)$, a~$\cal N$ liczbę niereszt w~tym przedziale.
Ponadto przez~$h$~oznaczmy
\begin{center}
\vspace*{-.1cm}
\color{magenta}rząd grupy klas ideałów pierścienia liczb całkowitych ciała
$\mathbb Q(\sqrt{-p})$.
\vspace*{-.1cm}
\end{center}
Zapewne nawet wśród Czytelników \textit{Delty} jest wielu takich, dla których
powyższa definicja jest niezrozumiała. Nie szkodzi -- dla nas istotne jest tylko
to, że liczba~$h$ jako rząd (czyli liczba elementów) grupy jest dodatnia.
Okazuje się bowiem, że liczba ${\cal R}-{\cal N}$ wynosi $h$ albo $3h$ w~zależności od
tego, czy $q\equiv 7\pmod{8}$, czy też  $q\equiv 3\pmod{8}$, a~stąd oczywiście ${\cal R}>{\cal N}$! Pierwsze
dowody wspomnianej równości dotyczącej $\cal R-\cal N$
podał Dirichlet
w~połowie XIX~wieku, tworząc w~ten sposób zręby analitycznej teorii liczb.
Natomiast pierwszy całkowicie elementarny (to nie znaczy, że prosty!) dowód
podał matematyk amator H.L.S.~Orde w~1978 roku w~artykule \href{https://mathoverflow.net/questions/25263/most-squares-in-the-first-half-interval}{\textit{On
Dirichlet's class number formula}} (Journal of London Mathematical Society 18). 
% https://mathoverflow.net/questions/25263/most-squares-in-the-first-half-interval
Ten dowód można
znaleźć w~5. rozdziale książki profesora Władysława Narkiewicza \textit{Classical
Problems in
Number Theory} (PWN, Warszawa 1986).
\eject

\normalsize




